\section{完整 Workflow 流程}

\subsection{六步流程}

\begin{enumerate}
    \item \textbf{Step 1：初始求解}——基于 Prompt 和 Gemini 2.5 Pro 产生一个初始解
    \item \textbf{Step 2：Self-Improvement}——共享 Step 1 的上下文，继续 refine（因为 Step 1 可能因 Thinking Budget 耗尽而仓促结束）
    \item \textbf{Step 3：Verification}——对 Step 2 的解进行验证，输出 verdict（有效/无效）和 bug report
    \item \textbf{Step 4：Bug Report}——提取具体的错误位置和原因
    \item \textbf{Step 5：Correction}——基于 bug report 修正 solution
    \item \textbf{Step 6：循环}——Step 3-5 不断循环
\end{enumerate}

\begin{importantbox}{终止条件}
\begin{itemize}
    \item \textbf{接受}：某个 solution 连续 \textbf{5 次}通过 verification $\rightarrow$ 认为解决了问题
    \item \textbf{拒绝}：循环 \textbf{10 次}仍未通过 $\rightarrow$ 放弃当前方向，重新 sample
\end{itemize}
\end{importantbox}

\subsection{Evaluator-Optimizer 模式}

\begin{knowledgebox}{与 Building Effective Agents 中的经典模式对应}
这个 Workflow 本质是 \textbf{Evaluator-Optimizer} 模式：
\begin{itemize}
    \item Generator（生成器）产生 solution
    \item Evaluator（评估器）验证并给出 feedback
    \item 接受则输出，拒绝则基于 feedback 重新生成
\end{itemize}
\end{knowledgebox}

\subsection{初始解质量的决定性作用}

\begin{warningbox}{成功的关键在于初始解}
如果 Step 2 生成的初始解与正确解有较大的 overlap，那么经过后续的验证-修正循环大概率能最终 solve 问题。但如果初始解的方向就是错的（偏差太远），后续的纠正大概率救不回来。
\end{warningbox}

\subsection{本章小结}

整个 Workflow 是一个标准的 Generate-Verify-Correct 迭代循环，终止条件基于连续通过/失败次数的计数。

